\section{Actual affine catalogue weight and cubic aggregation}
\label{sec:affine-catalogue-weight-overlap}

The ray-period argument reduces the remaining non-arm-level collinear term
to the Euler-totient weight of the divisor labels that actually arise from
canonical repeated kernels.  We now compute that weight and aggregate its
cubic energy without assuming that different kernel catalogues are
disjoint.

For a positive kernel triple
$d=(d_U,d_V,d_W)$ put $D(d)=d_Ud_Vd_W$ and
\[
 \mathcal D(d)=\{(e_U,e_V,e_W):e_Z\mid d_Z\},\qquad
 w(e)=\varphi(e_U)\varphi(e_V)\varphi(e_W).
\]
For a threshold $T$, define
\[
 L_T(d)=\sum_{\substack{e\in\mathcal D(d)\\e_Ue_Ve_W>T}}w(e).
\]

\begin{theorem}[Actual catalogue weight and finite tail]
\label{thm:affine-actual-catalogue-tail}
For every positive triple $d$,
\[
 \sum_{e\in\mathcal D(d)}w(e)=D(d)
\]
and
\[
 D(d)\le L_T(d)+T\tau(d_U)\tau(d_V)\tau(d_W).
\]
Moreover, if $T<D(d)$, then the top label gives
\[
 \varphi(d_U)\varphi(d_V)\varphi(d_W)\le L_T(d).
\]
\end{theorem}

\begin{proof}
The three sums separate, and
$\sum_{e\mid n}\varphi(e)=n$.  On the complementary range
$e_Ue_Ve_W\le T$, one has
$w(e)\le e_Ue_Ve_W\le T$, and there are at most
$\tau(d_U)\tau(d_V)\tau(d_W)$ labels.  Finally, $e=d$ occurs in the large
range whenever $T<D(d)$.
\end{proof}

At a canonical exceptional affine point let
$d=(K(U),K(V),K(W))$, where $K$ retains the exact prime powers of exponent
at least two.  With
\[
 M=\left\lfloor\frac{c^6}{4R}\right\rfloor,
 \qquad N=M-1,
\]
the established arm and excess inequalities give
\[
 27D(d)\le c^{20},\qquad
 D(d)>\frac{Rc^{14}}{8192},\qquad
 N^2<\frac{c^{12}}{16R^2}.
\]
The elementary bound $\tau(n)\ll_\epsilon n^\epsilon$, applied with
$\epsilon=1/40$ to the three components, therefore yields uniformly
\[
 \frac{D(d)-L_{N^2}(d)}{D(d)}
 \ll c^{-3/2}.
\]
Thus asymptotically all of the actual totient weight lies above the sharp
box-square determinant threshold.

Let $\mathcal K$ be the realized canonical kernel triples, and let
$m_\kappa$ be the size of the corresponding class in the exceptional-set
partition.  If $n(e)$ is the number of exceptional points whose canonical
kernel is divisible componentwise by $e$, then
\[
       n(e)=\sum_{\kappa:e\mid\kappa}m_\kappa.
\]

\begin{theorem}[Monotone catalogue aggregation]
\label{thm:affine-monotone-catalogue-aggregation}
For the determinant-large part of the third-gcd energy,
\[
 \sum_{\kappa\in\mathcal K}m_\kappa^3L_T(\kappa)
 \le
 \sum_{\substack{e\\e_Ue_Ve_W>T}}w(e)n(e)^3.
\]
Consequently, at $T=N^2$ the right side is at least
\[
 \bigl(1-O(c^{-3/2})\bigr)
 \sum_{\kappa\in\mathcal K}D(\kappa)m_\kappa^3.
\]
\end{theorem}

\begin{proof}
For each fixed label,
\[
 \sum_{\kappa:e\mid\kappa}m_\kappa^3
 \le
 \left(\sum_{\kappa:e\mid\kappa}m_\kappa\right)^3=n(e)^3.
\]
Multiply by $w(e)$, sum over $e$, and interchange the two finite sums.  The
asymptotic consequence follows from Theorem~\ref{thm:affine-actual-catalogue-tail}.
\end{proof}

This theorem controls energy, including its diagonal part.  It does not by
itself force two kernel classes to share a large label.  If every point
introduces a new catalogue, all relevant occupancies may still be one.  A
bound for deduplicated catalogue weight or singleton novelty remains
necessary for a genuine reuse conclusion.

The form best matched to ray capacities is shifted.  For nonnegative
weights and integer occupancies put
\[
 W=\sum_iw_i,\qquad I=\sum_iw_in_i,\qquad
 J=\sum_iw_i(n_i-1).
\]
Weighted H\"older gives
\[
 J^3\le W^2\sum_iw_i(n_i-1)^3.
\]
Since every occurring label has $n_i\ge1$, the termwise identity
$n_i=1+(n_i-1)$ gives $I=W+J$.  Thus the same estimate is the explicit
overlap inequality
\[
 (I-W)^3\le W^2\sum_iw_i(n_i-1)^3.
\]
Hence individual caps $(n_i-1)^3\le X_i$ imply
\[
                         J^3\le W^2\sum_iw_iX_i.
\]
Labels may use different supporting lines and different $X_i$; no raw
number-of-rays factor occurs.  This also avoids the factor four used when
the earlier paper converted shifted cubes to $n_i^3$.

The finite tail cannot be strengthened by simply subtracting $T$.  The
positive pairwise-coprime powerful triple
$(d_U,d_V,d_W)=(9,25,1)$ at $T=5$ has total weight $225$, small-label weight
$1+2+4=7$, and therefore
\[
                         L_5(d)=218<220=225-5.
\]
This is a complete counterexample only to $L_T(d)\ge D(d)-T$.  It leaves the
correct divisor-count tail and the affine route intact.

The Lean companion checks the exact catalogue identities, tail bounds,
top-label baseline, cubic aggregation, polynomial and shifted weighted
H\"older inequalities, and this precise counterexample.  The active next
steps are a signed period/capture theorem for every non-arm-level supporting
line, the two negative-slope arm-level caps, and an arithmetic bound on
deduplicated catalogue novelty.  No full-premise counterexample to those
corrected targets is known.
